%! AP Statistics — Unit 8, Lesson 8.5 — Guided Notes (Answer Key)
\documentclass[10pt]{article}
\usepackage{apstats-article}
\usepackage{apstats-key}
\newcommand{\TallMath}[1]{$\displaystyle #1\rule[-1.4em]{0pt}{3.2em}$}

\begin{document}

\pageheader{Unit 8, Lesson 8.5}{Guided Notes --- Answer Key}
\namedateperiod

\vspace{0.1in}

\begin{objectivebox}
By the end of this lesson, I will be able to:
\begin{itemize}
    \item \ansline{Distinguish between a chi-square test for homogeneity and a test for independence based on the study design.}
    \item \ansline{State appropriate $H_0$ and $H_a$ for each type of chi-square test.}
    \item \ansline{Verify the conditions required for chi-square tests for two-way tables.}
\end{itemize}
\end{objectivebox}

\vspace{0.08in}

\begin{vocabbox}
\termblanklong{test for homogeneity}
\ansline{A chi-square test that compares the distribution of one categorical variable across two or more populations or treatment groups.}
\termblanklong{test for independence}
\ansline{A chi-square test that determines whether two categorical variables are associated in a single population.}
\termblanklong{association}
\ansline{A relationship between two variables such that knowing one variable's value changes the probability distribution of the other.}
\termblanklong{stratified random sample}
\ansline{A sampling method in which the population is divided into subgroups (strata) and a random sample is drawn from each stratum.}
\end{vocabbox}

\vspace{0.08in}

\begin{notesbox}{Section 1: Two Study Designs}
\begin{minipage}[t]{0.47\linewidth}
    \textbf{Design Type 1:} One sample \\
    We measure \textbf{two} categorical variables on each individual.\\
    We want to know if \ans{there is an association} \\
    between the two variables.\\[6pt]
    $\Rightarrow$ This is a test for \ans{independence}.
\end{minipage}
\hfill
\begin{minipage}[t]{0.47\linewidth}
    \textbf{Design Type 2:} Multiple groups \\
    We measure \textbf{one} categorical variable on each group.\\
    We want to know if \ans{the distributions are the same} \\
    across the groups.\\[6pt]
    $\Rightarrow$ This is a test for \ans{homogeneity}.
\end{minipage}

\vspace{0.1in}
\textbf{Decision Questions:}
\begin{enumerate}
    \item Were data collected from \ans{one} sample(s) or \ans{multiple} groups?
    \item Were \ans{one} or \ans{two} categorical variables measured per individual?
\end{enumerate}
\end{notesbox}

\vspace{0.08in}

\begin{notesbox}{Section 2: Hypotheses}
\renewcommand{\arraystretch}{2.0}
\begin{tabularx}{\linewidth}{l X X}
    & \textbf{Test for Homogeneity} & \textbf{Test for Independence} \\
    \hline
    $H_0$ & There is \ans{no} difference in the distributions of \ans{[variable]} across \ans{[groups]}.
           & There is \ans{no} association between \ans{[var 1]} and \ans{[var 2]} in the population. \\
    $H_a$ & There \ans{is} a difference in distributions (at least one group differs).
           & There \ans{is} an association between the two variables. \\
\end{tabularx}
\end{notesbox}

\vspace{0.08in}

\begin{notesbox}{Section 3: Conditions}
\textbf{1. Independence condition}
\begin{itemize}
    \item For a test for homogeneity: data come from a \ans{stratified random sample} or a \ans{randomized experiment}.
    \item For a test for independence: data come from a \ans{simple random sample (SRS)}.
    \item When sampling without replacement: $n \le$ \ans{10}\% of $N$.
\end{itemize}
\textbf{2. Large-counts condition}
\begin{itemize}
    \item All \ans{expected} counts must be greater than \ans{5}.
    \item Formula: $E = \dfrac{(\text{row total})(\text{\ans{column total}})}{n}$
\end{itemize}
\end{notesbox}

\vspace{0.08in}

\begin{notesbox}{Section 4: Classify Each Scenario}
\begin{enumerate}[label=\alph*.]
    \item A researcher randomly selects 100 students from each of three high schools
          (Schools A, B, C) and asks which extracurricular activity (Sports/Arts/Neither)
          each student participates in. \\[4pt]
          Test type: \ans{test for homogeneity} \quad Because: \ansline{Three separate random samples (one per school) measure one categorical variable (extracurricular activity).}

    \item A random sample of 300 adults is asked their political affiliation
          (Democrat/Republican/Independent) and their view on a new policy (Support/Oppose).\\[4pt]
          Test type: \ans{test for independence} \quad Because: \ansline{One random sample measures two categorical variables (political affiliation and policy view) on each individual.}
\end{enumerate}
\end{notesbox}

\end{document}
